\chapter{IMPLEMENTATION AND RESULTS}
\label{ch:results}

\section{Development Environment}

\begin{table}[H]
\centering
\small
\caption{Tools}
\label{tab:tools}
\begin{tabularx}{\textwidth}{@{}lX@{}}
\toprule
\textbf{Tool} & \textbf{Use} \\
\midrule
STM32CubeMX 6.18 & peripheral, clock and middleware configuration of both projects \\
STM32CubeIDE 2.2 / GNU Arm toolchain 14.3 & firmware build \\
STM32CubeProgrammer CLI & flashing the NUCLEO-G431RB over SWD \\
OpenSTLinux (kernel 6.6), \texttt{remoteproc}, SSH over USB & loading and starting the M4 firmware on the DK1 \\
MSYS2 GCC and GNU Make & host unit tests \\
Git / GitHub & version control; each verified milestone is tagged \\
\bottomrule
\end{tabularx}
\end{table}

The M4 firmware is deployed by copying the ELF file to \texttt{/lib/firmware} and writing to the \texttt{remoteproc} sysfs interface (Appendix~\ref{app:procedures}).

\section{Unit Tests}

The hardware-independent modules are tested on a PC with a minimal self-written test harness (no external framework), a fake CAN driver that records every transmitted frame and can simulate a busy driver, and a fake millisecond clock that only advances when a test says so. Table~\ref{tab:unittests} summarises the suites.

\begin{table}[H]
\centering
\small
\caption{Unit test results}
\label{tab:unittests}
\begin{tabularx}{\textwidth}{@{}lccX@{}}
\toprule
\textbf{Suite} & \textbf{Tests} & \textbf{Checks} & \textbf{Scope} \\
\midrule
ISO-TP & 26 & 115 & SF/FF/CF/FC in both directions, padding, block size, STmin in ms and \textmu s encoding, N\_Bs and N\_Cr timeouts, FC.WAIT limit, FC.OVFLW, receive buffer overflow, wrong sequence number, SF interrupting a segmented reception, sequence number wrap-around (200 bytes), driver-busy retry, and a 4095-byte end-to-end transfer between two links on a simulated bus \\
E2E & 8 & 58 & CRC-8 catalogue value (\texttt{"123456789"} $\rightarrow$ 0x4B), counter wrap 15$\rightarrow$0, single-bit corruption, repeated frame, lost frames and resynchronisation, short frame \\
DTC manager & 11 & 35 & status bits after pass/fail, confirmation after consecutive failures with debounce reset, healing keeps history bits, operation-cycle rules for \textit{pendingDTC}, clear, read by status mask, balanced lock callbacks \\
UDS server & 11 & 69 & every service with positive and negative responses, suppress-positive-response bit, S3 session timeout, multi-DID read, response-too-long, reset only in the extended session and only after the response, functional-addressing NRC suppression \\
\midrule
\textbf{Total} & \textbf{56} & \textbf{277} & all pass, 0 failures \\
\bottomrule
\end{tabularx}
\end{table}

The 4095-byte end-to-end test also checks the exact frame count: one first frame, 585 consecutive frames and 74 flow-control frames for a block size of 8. An error in the expected count found while writing this test (584 instead of 585 consecutive frames, because $4089 = 584 \times 7 + 1$) illustrates why such arithmetic belongs in automated tests rather than in comments. Similarly, one DTC test initially expected a DTC with status 0x00 to be reported for the mask 0xFF; ISO~14229 reports a DTC only if \texttt{(status \& mask) != 0}, so the implementation was right and the expectation was corrected.

\section{Milestone M1: Sensor ECU in Loopback}

The Sensor ECU firmware was flashed with the STM32CubeProgrammer CLI and run in FDCAN internal loopback, where the controller receives its own frames. The pass criterion was: transmit rate as designed, every 0x100 frame received back, no dropped frame, no bus-off, and a stable ADC reading that follows the potentiometer. Figure~\ref{fig:m1} shows the log and Table~\ref{tab:m1} the results.

\reportimage{m1_sensor_loopback_uart.png}{0.95}{Sensor ECU log in FDCAN internal loopback (one line per second)}{fig:m1}

\begin{table}[H]
\centering
\small
\caption{Milestone M1 results}
\label{tab:m1}
\begin{tabularx}{\textwidth}{@{}lXc@{}}
\toprule
\textbf{Check} & \textbf{Observation} & \textbf{Result} \\
\midrule
Transmit rate & +110 frames per second (100$\times$0x100 + 10$\times$0x101) & \cmark \\
Loopback reception & \texttt{rx100} +100 per second; \texttt{rx = tx - 2} constant (two frames in flight when the log is printed) & \cmark \\
Endurance & more than 59\,000 frames, \texttt{drop=0}, \texttt{busoff=0} & \cmark \\
ADC & follows the potentiometer (0 to 104.8\,km/h observed during the sweep); $\pm$1--2\,LSB at rest & \cmark \\
Scaling & e.g. ADC 865 $\rightarrow$ $865 \times 2500/4095 = 528$ $\rightarrow$ 52.8\,km/h & \cmark \\
Fault injection (B1) & not yet exercised on hardware & open \\
\bottomrule
\end{tabularx}
\end{table}

\section{Milestone M2: Control ECU in Loopback}

In loopback the Control ECU also plays the Sensor ECU: its monitor task sends an E2E-protected 0x100 frame every 10\,ms with a speed ramp from 0 to 250\,km/h, so the whole chain from the FDCAN interrupt to the PWM output is exercised on the DK1 alone. Figure~\ref{fig:m2} shows the log and Table~\ref{tab:m2} the results.

\reportimage{m2_control_loopback_uart.png}{0.95}{Control ECU log on the STM32MP157 Cortex-M4 in FDCAN internal loopback}{fig:m2}

\begin{table}[H]
\centering
\small
\caption{Milestone M2 results}
\label{tab:m2}
\begin{tabularx}{\textwidth}{@{}lXc@{}}
\toprule
\textbf{Check} & \textbf{Observation} & \textbf{Result} \\
\midrule
Reception & \texttt{e2e ok} +100 per second, \texttt{rx = tx - 1} & \cmark \\
E2E & 30\,900 frames: 0 CRC, 0 lost, 0 repeated, 0 sequence errors & \cmark \\
Signal freshness & \texttt{age = 10\,ms} (timeout threshold 100\,ms) & \cmark \\
Output & PWM tracks speed, e.g. 193.4\,km/h $\rightarrow$ 77.3\,\% & \cmark \\
Faults & timeout active at start-up, cleared automatically once data arrives; no active fault afterwards & \cmark \\
RTOS & 0 queue overflows; 11\,032 bytes of the 16\,KB heap free & \cmark \\
\bottomrule
\end{tabularx}
\end{table}

\section{Milestone M3: UDS and DTC Management on the Control ECU}

The UDS server, the DTC manager and an ISO-TP link were integrated into the Control ECU, and the on-target self-test described in Section~\ref{sec:diag_design} was run in FDCAN internal loopback. Listing~\ref{lst:m3} shows the log; the complete log is stored in the repository (\texttt{docs/logs/m3\_uds\_selftest.log}). All 14 steps passed on the first hardware run.

\begin{lstlisting}[language={},numbers=none,caption={On-target UDS self-test on the Cortex-M4 (UART7 log, abridged)},label={lst:m3}]
[TEST]  1/14 TesterPresent                    PASS
[TEST]  2/14 ReadDID F195 (ECU sends FF+CF)   PASS
[TEST]  3/14 ReadDID x4 (ECU receives FF+CF)  PASS
[TEST]  4/14 Unknown service -> NRC 11        PASS
[TEST]  5/14 Unknown DID -> NRC 31            PASS
[TEST]  6/14 Reset in default -> NRC 7E       PASS
[TEST]  7/14 Session extended                 PASS
[TEST]  8/14 Confirmed DTCs = 0               PASS
[CTRL] fault OUTPUT_LOCAL  SET
[TEST]  9/14 Confirmed DTCs = 1               PASS
[TEST] 10/14 DTC list B1A00-11 0x2F           PASS
[CTRL] fault OUTPUT_LOCAL  cleared
[TEST] 11/14 Clear all DTCs                   PASS
[TEST] 12/14 Confirmed DTCs = 0 again         PASS
[CTRL] soft reset: new operation cycle
[TEST] 13/14 Soft reset                       PASS
[TEST] 14/14 Session default                  PASS
[TEST] UDS self-test finished: 14/14 passed
[CTRL] speed=249.5 km/h age=10ms pwm=99.8% | e2e ok=499 lost=0 crc=0 rep=0
       seq=0 | faults=0x00 dtc=0 | rx=537 tx=538 busoff=0 qovf=0 uds=38 heap=9744
\end{lstlisting}

\begin{table}[H]
\centering
\small
\caption{Milestone M3 results}
\label{tab:m3}
\begin{tabularx}{\textwidth}{@{}lXc@{}}
\toprule
\textbf{Check} & \textbf{Observation} & \textbf{Result} \\
\midrule
Segmented transfers & 13-byte response from the ECU and 9-byte request to the ECU, both with FF/CF and flow control & \cmark \\
Negative responses & NRC 0x11 (unknown service), 0x31 (unknown DID), 0x7E (reset in default session) & \cmark \\
Fault to DTC & local fault injected $\rightarrow$ DTC B1A00-11 confirmed with status 0x2F (TF, TFTOC, PDTC, CDTC, TFSLC) & \cmark \\
Clear & 0x14 FFFFFF $\rightarrow$ no confirmed DTC & \cmark \\
Sessions and soft reset & extended session with P2/P2*, soft reset starts a new operation cycle & \cmark \\
No side effects & 38 diagnostic frames exchanged while sensor data kept flowing: 0 E2E errors, 0 queue overflows; 9.7\,KB heap free & \cmark \\
Hard reset (0x11 01) & cannot be exercised by a tester running on the same core & open \\
\bottomrule
\end{tabularx}
\end{table}

\section{Milestone M4: Two-Node Bus with Fault Injection}

\subsection{Wiring}

Table~\ref{tab:wiring} lists the verified connections. Each transceiver module carries a 120\,$\Omega$ terminator, so the bus measures about 60\,$\Omega$ between CANH and CANL; both boards share a common ground.

\begin{table}[H]
\centering
\small
\caption{Verified bus wiring}
\label{tab:wiring}
\begin{tabularx}{\textwidth}{@{}lXlX@{}}
\toprule
\textbf{Signal} & \textbf{NUCLEO-G431RB} & \textbf{SN65HVD230} & \textbf{STM32MP157D-DK1} \\
\midrule
FDCAN1\_TX & PA12, CN10 pin 12 & CTX & PA12, CN13 pin 9 (Arduino D14) \\
FDCAN1\_RX & PA11, CN10 pin 14 & CRX & PA11, CN13 pin 10 (Arduino D15) \\
Supply & 3V3 / GND (CN6) & 3V3 / GND & CN16 pin 4 / pin 6 \\
\bottomrule
\end{tabularx}
\end{table}

\subsection{Diagnosing the bus bring-up}

The first power-up on the bus did not work, and the fault was located without an oscilloscope by reading the FDCAN status registers of both nodes (the Nucleo through the ST-LINK in hot-plug mode, the DK1 from Linux). The last error code (LEC) of the protocol status register was the key observable:

\begin{enumerate}[nosep]
  \item The Sensor ECU reported LEC\,=\,3 (\textit{acknowledge error}) with its transmit error counter at 128. A node only reaches the acknowledge slot if it has read back every bit it sent, so the path Nucleo\,$\rightarrow$\,transceiver\,$\rightarrow$\,bus\,$\rightarrow$\,transceiver\,$\rightarrow$\,Nucleo was proven; only an acknowledging receiver was missing.
  \item On the DK1 the receive pin PA11, sampled 50\,000 times per second from Linux, never went low while the Nucleo was retransmitting: no signal reached the DK1.
  \item Swapping the two transceiver modules on the Nucleo reproduced LEC\,=\,3 with the second module as well, so both modules were good and the fault was on the DK1 side. The DK1 wires were moved to the documented header pins (CN13 pins 9 and 10).
  \item A cyclic status frame (0x200) was added to the Control ECU so that it also transmits. The DK1 then reported LEC\,=\,3 too: each node read back its own frames, but neither received the other, which located the fault in the link between the two modules; the CANH/CANL wires had been removed during module testing.
  \item After reconnecting, both nodes reported LEC\,=\,5 (\textit{bit-0 error}) and cycled through bus-off. To separate a polarity error from other causes, the DK1 controller was switched to bus-monitoring (listen-only) mode from Linux: it received the Sensor ECU frames with valid E2E checks, which proved the polarity correct. With stable connections the bus then ran without errors.
\end{enumerate}

During this phase the Sensor ECU went through several hundred bus-off events and recovered from each automatically, which exercises requirement R8 on real hardware.

\subsection{Fault injection test}

Both logs were captured simultaneously for 65\,s while the potentiometer was swept and button B1 of the Sensor ECU was used to inject faults. Table~\ref{tab:m4} summarises the run; the complete logs are in \texttt{docs/logs/}.

\begin{table}[H]
\centering
\small
\caption{Milestone M4: fault injection over the physical bus}
\label{tab:m4}
\begin{tabularx}{\textwidth}{@{}lXX@{}}
\toprule
\textbf{Time} & \textbf{Sensor ECU} & \textbf{Control ECU} \\
\midrule
23--33\,s & potentiometer sweep, 142--233\,km/h & received speed and PWM follow \\
39.7\,s & B1: SILENT, frame 0x100 stopped & 39.8\,s: SPEED\_TIMEOUT set, PWM 0\,\%, DTC U0100-87 confirmed \\
46.7\,s & B1: OUT\_OF\_RANGE, 300.0\,km/h sent & timeout cleared, SPEED\_RANGE set, PWM 0\,\%, DTC P0501-00 confirmed \\
55.6\,s & B1: back to normal & SPEED\_RANGE cleared, PWM 64.5\,\% again; confirmed DTCs retained \\
whole run & 110 frames/s, no frame dropped & 100 frames/s, 0 CRC, lost or repeated errors, 0 bus-off, 0 queue overflows \\
\bottomrule
\end{tabularx}
\end{table}

The timeout was reported one log line (100\,ms resolution) after the Sensor ECU changed mode, consistent with the 100\,ms timeout threshold. The DTC count stayed at three after all faults had healed: confirmed DTCs are kept until a tester clears them, as ISO~14229 requires. One of the three, U0400-81, was stored when the bus first came up: the Sensor ECU flushed three stale frames from its transmit FIFO with old alive counters, and the E2E check correctly classified them as out of sequence.

\section{Bring-up of the Control ECU}
\label{sec:bringup}

The first run of the Control ECU on the DK1 did nothing observable. No debugger connection to the M4 was available in production mode (the A7 owns the M4 reset), and the UART log was not yet readable. The firmware state was therefore inspected from Linux, which can read the M4 SRAM (aliased at 0x3002\,0000) and all peripheral registers through \texttt{/dev/mem}; symbol addresses and structure offsets were taken from the ELF file with \texttt{nm} and \texttt{gdb}. Two independent defects were found.

\subsection{Issue 1: missing SysTick handler}

\textbf{Observation.} \texttt{xSchedulerRunning = 1} and the current task was the idle task, but the FreeRTOS tick count \texttt{xTickCount} stayed at 0 and the HAL tick \texttt{uwTick} was frozen at 19. All application tasks were blocked waiting for a tick that never came.

\textbf{Root cause.} STM32CubeMX generated \texttt{USE\_CUSTOM\_SYSTICK\_HANDLER\_IMPLEMENTATION = 1}, which tells the CMSIS-RTOS2 layer not to define \texttt{SysTick\_Handler} because the interrupt file will. With TIM7 selected as HAL time base, however, the interrupt file contained no such handler. The symbol table showed \texttt{SysTick\_Handler} as a weak alias of \texttt{Default\_Handler}, an endless loop: the first RTOS tick trapped the CPU at the SysTick priority, which also blocked the TIM7 interrupt and froze \texttt{uwTick}.

\textbf{Fix.} The flag is overridden to 0 inside a \texttt{USER CODE} section of \texttt{FreeRTOSConfig.h}, which STM32CubeMX preserves on regeneration, so that \texttt{cmsis\_os2.c} provides the standard handler. After the fix the symbol is a strong definition and the tick advances at 1\,kHz.

\subsection{Issue 2: FDCAN kernel clock stuck in a glitch-free multiplexer}

\textbf{Observation.} After the first fix the tasks ran, but \texttt{CCCR.INIT} remained 1, no frame was transmitted and the HAL state was \texttt{BUSY} with no error code, meaning that \texttt{HAL\_FDCAN\_Start()} had succeeded. Clearing \texttt{INIT} directly from Linux had no effect either, which pointed to a missing protocol clock rather than to the firmware.

\textbf{Investigation.} Register reads after a fresh Linux boot showed the state left by Linux (Table~\ref{tab:clockstate}). The clock calibration unit of FDCAN1 was in bypass (\texttt{CCFG.BCC = 1}) and therefore not involved.

\begin{table}[H]
\centering
\small
\caption{FDCAN clock state left by Linux after boot}
\label{tab:clockstate}
\begin{tabularx}{\textwidth}{@{}llX@{}}
\toprule
\textbf{Register field} & \textbf{Value} & \textbf{Meaning} \\
\midrule
\texttt{RCC\_FDCANCKSELR} & 3 & FDCAN kernel clock taken from PLL4\_R \\
\texttt{RCC\_PLL4CR.DIVREN} & 0 & PLL4\_R output disabled (unused by Linux) \\
\texttt{RCC\_OCENSETR.HSEKERON} & 0 & HSE kernel clock output not enabled \\
\texttt{RCC\_TZCR.TZEN} & 1 & RCC in secure mode; oscillator \textit{disable} register is secure-only \\
\bottomrule
\end{tabularx}
\end{table}

\textbf{Root cause.} STM32 kernel clock multiplexers are glitch-free: a switch completes only while both the current and the new input are running~\cite{RM0436}. The firmware switched the multiplexer from PLL4\_R to HSE while PLL4\_R was disabled, so the switch never completed and the FDCAN core was left without any clock. The hypothesis was confirmed experimentally: enabling PLL4\_R for a few milliseconds from Linux immediately let \texttt{INIT} clear and frames flow; PLL4\_R was then disabled again and FDCAN kept running from HSE.

\textbf{Fix.} \texttt{CanDrv\_PrepareKernelClock()}, called before the HAL initialises FDCAN, sets the HSE kernel clock enable, temporarily enables the output of the currently selected PLL if it is off, switches to HSE, waits 1\,ms, restores the PLL output exactly as Linux left it, and reads back the effective kernel frequency. The fix was verified after a fresh reboot with no manual register access: the firmware alone brought FDCAN up at 24\,MHz.

\subsection{Other observations}

\begin{itemize}
  \item \textbf{FreeRTOS heap.} The default 3\,KB heap was too small for five tasks (three application tasks plus the idle and timer-service tasks) with newlib re-entrancy enabled; it was raised to 16\,KB, of which 11\,KB remain free.
  \item \textbf{DLC encoding.} The STM32MP1 HAL encodes the data length in bits 19:16 of the header word, whereas the STM32G4 HAL uses the plain byte count. A compile-time assertion in each driver guards against silently using the wrong encoding.
  \item \textbf{HAL tick rate.} In production mode the HAL tick driven by TIM7 runs at about 2\,kHz, because the timer clock configured by Linux differs from the value assumed by STM32CubeMX. The application uses the FreeRTOS tick, which is correct; the HAL tick is an open point.
\end{itemize}

\section{Resource Usage}

\begin{table}[H]
\centering
\small
\caption{Memory footprint (Debug builds)}
\label{tab:memory}
\begin{tabular}{@{}lrrl@{}}
\toprule
\textbf{Firmware} & \textbf{Code} & \textbf{RAM (data+bss)} & \textbf{Available} \\
\midrule
Sensor ECU (STM32G431) & 31.6\,KB & 2.5\,KB & 128\,KB flash / 32\,KB RAM \\
Control ECU (M4, \texttt{-O0}, with UDS/DTC) & 76.5\,KB & 24.8\,KB & 128\,KB code SRAM / 128\,KB data SRAM \\
\bottomrule
\end{tabular}
\end{table}
